\documentclass[journal,twoside,web]{ieeecolor2}
\usepackage{generic}
\usepackage{cite}
\usepackage{amsmath,amssymb,amsfonts}
\usepackage{algorithmic}
\usepackage{graphicx}
\usepackage{textcomp}
\usepackage{booktabs}
\usepackage{multirow}
\usepackage{array}
\usepackage{siunitx}
\usepackage{url}

\def\BibTeX{{\rm B\kern-.05em{\sc i\kern-.025em b}\kern-.08em
    T\kern-.1667em\lower.7ex\hbox{E}\kern-.125emX}}
\markboth{IEEE TRANSACTIONS ON BIOMEDICAL ENGINEERING}
{Huang Tan \MakeLowercase{\textit{et al.}}: Validating Riemannian Alignment for EEG Vigilance}
\begin{document}

\title{Validating Riemannian Parallel Transport Alignment for Cross-Subject EEG Vigilance Estimation}

\author{{Huang Tan,Xiangzhu Li,Li Zhang,Guangqiang Yin}
\thanks{Corresponding author: Guangqiang Yin}
\thanks{Huang Tan is with University of Electronic Science and Technology of China, Chengdu, 611731, China. (e-mail: tanhuang@std.uestc.edu.cn).}
\thanks{Xiangzhu Li and Li Zhang are with Shenzhen Clinical College of Integrated Chinese and Western Medicine, Guangzhou University of Chinese Medicine,No.528 Xinsha Road, Shajing Street, Shajing Subdistric, Bao'an District, Shenzhen City, GuangdongProvince, 518104 , China,  (e-mail: thyysjwk1bq@163.com,lxz142537@163.com).}
\thanks{Guangqiang Yin is with the School of Information and Software Engineering, University of Electronic Science and Technology of China, Chengdu 610057, China ,also with the Shenzhen Institute for Advanced Study,University of Electronic Science and Technology of China, Shenzhen 518110, China.(e-mail: yingq@uestc.edu.cn).}
}

\maketitle

\begin{abstract}
Cross-subject generalization remains a critical bottleneck for EEG-based
vigilance estimation, because subject-specific physiological factors
displace the SPD covariance matrices of different subjects to different
regions of the Riemannian manifold, degrading tangent space regression.
Existing Riemannian alignment methods have been validated on classification
tasks, but their effectiveness for continuous vigilance regression has not
been established. We apply Riemannian parallel transport alignment that
projects each subject's SPD matrices to a common reference at the identity,
removing subject-specific bias while preserving vigilance-related
variability. On the SEED-VIG dataset (23 subjects) under leave-one-subject-out
cross-validation, the alignment raises the mean Pearson correlation from
$0.585$ to $0.692$ ($p = 0.0007$, Cohen's $d = 0.824$), improving 21 of 23
subjects individually. The benefit persists on the 4-channel forehead setup
(retaining 98.3\% of full-setup performance) and matches or exceeds existing
Riemannian transfer-learning baselines. These results establish Riemannian
parallel transport as a practical, statistically robust pre-processing step
for cross-subject EEG vigilance regression.
\end{abstract}

\begin{IEEEkeywords}
EEG, vigilance estimation, Riemannian geometry, domain adaptation,
parallel transport, SPD manifold, cross-subject generalization, PERCLOS.
\end{IEEEkeywords}

\section{Introduction}
\label{sec:intro}

\IEEEPARstart{V}{igilance} monitoring from scalp electroencephalography
(EEG) is a long-standing goal of brain-computer interface (BCI)
research, with direct applications in driver fatigue detection
\cite{gao2024, cui2023},
neurological disorder diagnosis, and anesthesia depth monitoring.
In the SEED-VIG benchmark \cite{zheng2017}, twenty-three
subjects completed a 120-minute sustained-attention task while 17-channel
EEG and PERCLOS labels were recorded synchronously, providing a standard
test bed for cross-subject vigilance regression.

The dominant pipeline for EEG-based vigilance regression has three
stages: a filter bank that decomposes the signal into canonical frequency
bands ($\delta$, $\theta$, $\alpha$, $\beta$, $\gamma$); symmetric positive
definite (SPD) covariance features computed within each band; and a
regressor trained on the tangent space projection of those covariance
features. This Filter-Band Tangent Space (FBTS) pipeline has been shown
to be competitive for within-subject prediction
\cite{zheng2017, zanini2018, Rodrigues2019}. Recent
comprehensive reviews of EEG classification algorithms
\cite{lotte2018review} and EEG transfer-learning advances
\cite{wu2022tnsre} underline the importance of domain-adaptation
strategies for cross-subject generalization.

The remaining bottleneck is cross-subject generalization, which is
evaluated by leave-one-subject-out (LOSO) cross-validation. In the LOSO
protocol, the model is trained on $N-1$ subjects and tested on the held-out
one. The test subject's EEG covaries with the training subjects only
weakly, because of skull conductivity, electrode impedance, baseline
EEG power, and other physiological factors. These factors shift the
SPD covariance matrices of different subjects to different regions of
the SPD manifold. The FBTS pipeline projects the matrices to a tangent
space at a fixed reference point and trains a linear regressor on the
resulting vector. When the test subject's SPD matrices sit far from
the training reference, the tangent space projection is ill-conditioned
and the regressor generalizes poorly.

A natural remedy is to remove the subject-specific bias by aligning
the SPD distributions before projection. In the BCI literature, two
aligners are widely used. Euclidean alignment
\cite{he2018euclidean} whitens each subject by the inverse square
root of the subject's Euclidean mean covariance:
\begin{equation}
C'_{si} = M_s^{-1/2}\, C_{si}\, M_s^{-1/2},
\quad M_s = \frac{1}{T_s}\sum_i C_{si}.
\label{eq:euclidean}
\end{equation}
Riemannian alignment
\cite{barachant2010lvaica, barachant2012tbme, zanini2018, Rodrigues2019, zhuo2024} replaces
$M_s$ with the subject's Riemannian (Fr\'echet) mean $G_s$, giving
\begin{equation}
C'_{si} = G_s^{-1/2}\, C_{si}\, G_s^{-1/2},
\quad G_s = \arg\min_{G}\sum_i \delta_R^2(G, C_{si}),
\label{eq:riemann}
\end{equation}
where $\delta_R$ is the affine-invariant Riemannian distance on the SPD
manifold. Equation~\eqref{eq:riemann} is a parallel transport of the
subject's SPD matrices to a common reference at the identity. Geometric
parallel transport \cite{yair2019ptspd} is preferred over Euclidean whitening because it
preserves the affine-invariant structure of the SPD manifold \cite{moakher2005,pennec2006,congedo2017review} and
constitutes the right notion of centering on a curved space.

Riemannian alignment has been shown to improve BCI motor imagery
cross-session classification \cite{zanini2018, Rodrigues2019}, and recent
work has explored deep optimal transport on SPD manifolds for domain
adaptation \cite{ju2025}. However, it remains unclear whether
Riemannian alignment benefits continuous vigilance regression, how it
interacts with the FBTS pipeline under different channel configurations
and covariance estimators, and whether per-band alignment outperforms
global alignment for vigilance features.

\subsubsection{Contributions} In this paper we apply Riemannian parallel
transport alignment to cross-subject EEG vigilance regression on
SEED-VIG. Our contributions are fourfold. (i) On 23 subjects under
LOSO, Riemannian alignment raises mean COR from $0.585 \pm 0.230$ to
$0.692 \pm 0.206$ ($p = 0.0007$, $d = 0.824$, 21 of 23 subjects
improved) on the full 17-channel setup. (ii) The benefit extends to
the practical 4-channel forehead setup (permutation test
$p = 0.0001$, $d = 0.850$), retaining 98.3\%
of the full-setup performance. (iii) Ablations confirm that
the gain is attributable to the Riemannian geometry rather than
individual pipeline choices. (iv) The proposed approach is computationally
efficient and compatible with standard wearable EEG setups.

The remainder of the paper is organized as follows. Section~\ref{sec:method}
describes the FBTS pipeline, the SPD geometry, the three aligners, the
LOSO protocol, and the regression head. Section~\ref{sec:exp} reports
the main LOSO experiment, the few-shot calibration curve, the channel
configuration comparison, the ablations, and the frequency-band importance
analysis.
Section~\ref{sec:discussion} interprets the results and discusses
limitations, and Section~\ref{sec:conclusion} concludes.

\section{Methods}
\label{sec:method}

\subsection{Task and data}

We use the SEED-VIG dataset \cite{zheng2017}, which contains 23
subjects (sub01--sub23). Each subject completed a sustained-attention
driving task while 17-channel EEG and synchronized PERCLOS labels were
recorded. We use the standard 885-epoch segmentation provided by
SEED-VIG, with each epoch corresponding to 8 seconds of EEG. PERCLOS
values are continuous in $[0, 1]$. Performance is measured by Pearson
correlation (COR) and root mean square error (RMSE) between the
predicted and the ground-truth PERCLOS.

\subsection{Filter bank and SPD covariance features}

Following \cite{zheng2017}, the EEG signal is decomposed by a
fifth-order Butterworth filter bank into five canonical bands:
$\delta$ (1--4 Hz), $\theta$ (4--8 Hz), $\alpha$ (8--14 Hz), $\beta$
(14--31 Hz), and $\gamma$ (31--50 Hz). For each band and each epoch,
the $C \times C$ SPD covariance matrix is estimated by the Oracle
Approximating Shrinkage (OAS) estimator \cite{chen2010oas}, with
Ledoit-Wolf (LWF) and Sample Covariance Matrix (SCM) used in the
covariance-estimator ablation. The output of this stage is a
5-tuple of $C \times C$ SPD matrices per epoch, with $C \in \{17, 6, 4\}$
for the full, temporal, and forehead channel sets.

\subsection{Riemannian alignment of SPD matrices}

The three aligners we compare are defined by the choice of mean used
in \eqref{eq:euclidean} or \eqref{eq:riemann}:
\begin{itemize}
\item \textbf{NoAlignment} (baseline): $C'_{si} = C_{si}$, the
identity transform.
\item \textbf{EuclideanAlignment}: $C'_{si} = M_s^{-1/2} C_{si}
M_s^{-1/2}$, where $M_s$ is the subject's arithmetic mean of
$\{C_{si}\}$.
\item \textbf{RiemannianAlignment}: $C'_{si} = G_s^{-1/2} C_{si}
G_s^{-1/2}$, where $G_s$ is the subject's Riemannian (Fr\'echet) mean,
computed by the iterative algorithm in \cite{barachant2010lvaica, barachant2012tbme}.
\end{itemize}
Both Euclidean and Riemannian alignment are applied per band, because
different frequency bands exhibit different inter-subject variability
patterns. The alignment is fit on each subject's own training SPD
matrices only; the test subject is aligned with its own mean computed
from the test epochs, to avoid leaking test labels.

Geometrically, Equations~\eqref{eq:euclidean} and \eqref{eq:riemann}
are parallel transport maps that take the subject's SPD distribution
to a common reference at the identity. The two maps differ in their
notion of mean: $M_s$ is the Euclidean mean, which is a poor summary
on the curved SPD manifold, whereas $G_s$ is the Riemannian mean,
which is the correct notion of center on a curved manifold. The two
transforms are different in general; we refer to the difference as
the geometry of the mean. Section~\ref{sec:exp} reports that the
Riemannian mean produces strictly better regression performance.

\noindent\textbf{Relation to Zanini et al.}
The Zanini et al.\ alignment \cite{zanini2018} and the proposed method
both use Riemannian parallel transport and differ only in the reference
point (grand mean vs.\ identity). The affine invariance of the AIRM
distance guarantees that the Riemannian geometry is preserved under any
choice of reference. Empirically, the two reference points yield very
similar regression performance (see Section~\ref{sec:ablation},
Reference point). A pragmatic advantage of the identity reference is
that it circumvents the need to estimate a global Riemannian mean,
which reduces the per-subject computational overhead by approximately
$0.01$~s.

\subsection{Tangent space projection and regression}

For each band, the aligned SPD matrices are projected to the tangent
space at the identity using the matrix logarithm. The vectorized
upper-triangular entries of the logarithms from the five bands are
concatenated into a single feature vector. We then apply
\texttt{SelectKBest(f\_regression, k=150)} to remove low-importance
features, \texttt{StandardScaler} to normalize, and \texttt{RidgeCV}
regression to predict PERCLOS. To prevent data leakage, both
\texttt{SelectKBest} and \texttt{StandardScaler} are fit exclusively
on the training subjects of each LOSO fold and then applied to the
held-out test subject. Predictions are smoothed
temporally with a moving average of window 3 epochs.

The hyperparameters were selected through a nested cross-validation pilot
study on a subset of subjects: $k=150$ was chosen to balance the
trade-off between feature richness and overfitting risk (the full
tangent-space feature vector consists of $5 \times \binom{17}{2} + 17
= 765$ entries for the 17-channel configuration); \texttt{RidgeCV}
automatically selects the regularization strength $\alpha$ from a
logarithmic grid $\{0.01, 0.1, 1, 10, 100\}$ via internal leave-one-out
cross-validation on the training set; and the window size of 3
epochs was adopted as a lightweight temporal smoothing strategy
following prior EEG vigilance estimation pipelines
\cite{zheng2017,ang2012fbcsp}.

\subsection{Experimental protocol}

\subsubsection{Leave-one-subject-out (LOSO)} For each subject $s$ in
$\{1, \ldots, 23\}$, the model is trained on the other 22 subjects and
tested on subject $s$. Each fold uses the aligner, channel set, and
regressor specified by the experiment. The reported mean COR and RMSE
are averaged across the 23 folds. Statistical significance is
evaluated by a paired $t$-test on the per-subject COR values between
two aligners; Cohen's $d$ is reported as the effect size.

\subsubsection{Few-shot calibration} In the main LOSO experiment,
the test subject's Riemannian mean is estimated from all of its own
epochs (full alignment). In the few-shot experiment, we restrict the
number $n$ of epochs used to estimate $G_s$: with $n=0$, the test
subject is aligned using the grand mean of the training subjects'
Riemannian means; with $n \ge 1$, the $n$ calibration epochs from the
test subject are used to estimate $G_s$ and to align both the
calibration and the remaining test epochs. The reported COR is averaged
over five random seeds for each $n$.

\subsubsection{Channel configurations} We evaluate three channel sets:
all 17 channels, the 6 temporal channels (T7, T8, TP7, TP8, P7, P8),
and the 4 forehead channels (Fp1, Fp2, Fz, F1). The forehead set is of
practical interest for wearable EEG headsets that often place
electrodes on the forehead.

\subsubsection{Ethics statement} The SEED-VIG dataset
\cite{zheng2017} is publicly released for non-commercial research use.
The original data collection was conducted by the BCMI Laboratory at
Shanghai Jiao Tong University and was approved by the Institutional
Review Board of Shanghai Jiao Tong University. Written informed consent
was obtained from all participants prior to the experiment, in
accordance with the Declaration of Helsinki. No new data were collected
for this study. The present analysis is a secondary use of fully
de-identified recordings, and no further ethical approval was required
for the computational work reported here.

\section{Experiments and Results}
\label{sec:exp}

We compare three aligners (None, Euclidean, Riemannian) on three
channel sets (17ch, 6ch, 4ch). All experiments use the FBTS pipeline
described in Section~\ref{sec:method}. The full LOSO main result uses
23 subjects; the ablations use 10 subjects for runtime; the few-shot
experiment uses 5 random seeds per calibration size. Standard deviations
are reported across subjects unless otherwise noted.

\subsection{Main LOSO result (Table~\ref{tab:main})}

Table~\ref{tab:main} reports mean COR and RMSE for the 3$\times$3 grid
of aligners and channel sets on the full 23 subjects. Riemannian
alignment improves COR over the no-alignment baseline on all three
channel sets, with paired $t$-test $p=0.0007$ ($d=0.824$) on the full
17-channel setup. Euclidean alignment provides an intermediate
improvement on the 17-channel setup but is statistically
indistinguishable from baseline on the 6-channel temporal setup.

\begin{table}[!t]
\caption{LOSO Cross-Subject Vigilance Regression on SEED-VIG
(23 subjects). Mean $\pm$ std Pearson correlation (COR) and root mean
square error (RMSE) across the 23 LOSO folds. $p$ and Cohen's $d$ are
computed for the paired comparison of Riemannian vs.\ None on the same
channel set. For the 4-channel forehead setup, the reported $p$ is
from a Wilcoxon signed-rank test, with a permutation test
($n=10{,}000$ permutations) confirming the significance
($p_{\mathrm{perm}}=0.0001$).}
\label{tab:main}
\centering
\footnotesize
\setlength{\tabcolsep}{3pt}
\begin{tabular}{llcccc}
\toprule
\textbf{Channels} & \textbf{Aligner} & \textbf{COR} & \textbf{RMSE} & $\boldsymbol{p}$ & $\boldsymbol{d}$ \\
\midrule
\multirow{3}{*}{17ch (all)}
  & None       & $0.585\pm0.230$ & $0.278$ & ---    & ---   \\
  & Euclidean  & $0.655\pm0.218$ & $0.224$ & $0.005$ & $0.650$ \\
  & \textbf{Riemannian} & $\mathbf{0.692\pm0.206}$ & $\mathbf{0.195}$ & $\mathbf{0.0007}$ & $\mathbf{0.824}$ \\
\midrule
\multirow{3}{*}{6ch (temporal)}
  & None       & $0.659\pm0.185$ & $0.248$ & ---    & ---   \\
  & Euclidean  & $0.622\pm0.257$ & $0.214$ & n.s.   & n.s.  \\
  & \textbf{Riemannian} & $\mathbf{0.689\pm0.190}$ & $\mathbf{0.194}$ & n.s. & n.s. \\
\midrule
\multirow{3}{*}{4ch (forehead)}
  & None       & $0.626\pm0.196$ & $0.243$ & ---    & ---   \\
  & Euclidean  & $0.642\pm0.204$ & $0.208$ & n.s.   & n.s.  \\
  & \textbf{Riemannian} & $\mathbf{0.681\pm0.185}$ & $\mathbf{0.196}$ & $<0.0001^\dagger$ & $\mathbf{0.850}$ \\
\bottomrule
\end{tabular}
\end{table}

The per-subject COR values for the 17-channel setup are shown in
Fig.~\ref{fig:individual}. All three aligners are displayed; Riemannian
alignment improves 21 of the 23 subjects individually over the
no-alignment baseline. The two subjects with a small decrease
(sub23, sub02) are marked by red triangles; sub23 was the
lowest-baseline subject (rank 1) and sub02 was mid-range
(rank 8), both remaining below their baseline COR after
alignment. The 4-channel forehead setup
retains 98.3\% of the 17-channel Riemannian performance, which is
important for wearable EEG where full coverage is impractical.

\begin{figure*}[!t]
\centering
\includegraphics[width=\textwidth]{figures/individual_cor.pdf}
\caption{Per-subject LOSO Pearson correlation on the 17-channel setup,
sorted by baseline (None) COR. Three aligners are compared: None (gray),
Euclidean (light blue), and Riemannian (blue). Dashed horizontal lines
indicate the group mean under the None and Riemannian aligners, with the
mean values labeled at the right margin. Red triangles mark the two
subjects for which Riemannian alignment decreases COR relative to
baseline. The bottom-right annotation reports the number of improved
subjects (21/23), the paired $t$-test $p$-value, and Cohen's $d$.}
\label{fig:individual}
\end{figure*}

\subsection{Few-shot calibration curve}

Fig.~\ref{fig:fewshot} reports the LOSO correlation as a function of
the number of calibration epochs drawn from the test subject. With
$n=0$ calibration epochs the test subject is aligned using the
training subjects' grand Riemannian mean and the mean COR is $0.645$,
substantially above the no-alignment baseline ($0.585$) but below the
full-alignment performance ($0.692$). With $n=1$ calibration epoch the
mean COR increases to $0.681$ ($98.4\%$ of full performance). As
highlighted by the red star at $n=5$, five calibration epochs recover
the full-alignment performance ($0.693$). The curve remains stable
from $n=5$ to $n=100$ calibration epochs, with variance decreasing
as expected. This finding confirms that the proposed alignment protocol
does not require extensive calibration data from the target subject;
even a single calibration epoch recovers most of the benefit.

\begin{figure}[!t]
\centering
\includegraphics[width=\columnwidth]{figures/fewshot_curve.pdf}
\caption{Few-shot calibration curve: mean LOSO COR as a function of the
number of calibration epochs from the test subject. Error bars are
standard deviation across five random seeds. The dashed horizontal lines
indicate the full-alignment (blue) and no-alignment (gray) reference
levels. The red star marks $n=5$ calibration epochs, with the annotation
showing the corresponding COR and the percentage of full-alignment
performance recovered.}
\label{fig:fewshot}
\end{figure}

\subsection{Channel configuration comparison}

Fig.~\ref{fig:channels} compares the three aligners across the three
channel sets. Riemannian alignment is the best aligner on every
channel set. The absolute COR gain over baseline is $+0.107$ on the
17-channel full setup (paired $t$-test $p=0.0007$, Cohen's $d=0.824$,
marked ${}^{***}$ in the figure)
 and $+0.055$ on the 4-channel forehead setup (permutation test
$p=0.0001$, $d=0.850$). On the 6-channel temporal setup, the
no-alignment baseline is already relatively strong ($0.659$), and the
Riemannian improvement ($+0.030$) is smaller and not statistically
significant (marked n.s.). Importantly, the 4-channel forehead setup
retains 98.3\% of the full 17-channel Riemannian performance
($0.681$ vs $0.692$), as annotated in the figure, demonstrating the
practical viability of Riemannian alignment on wearable EEG systems
with limited channel count.

\begin{figure}[!t]
\centering
\includegraphics[width=\columnwidth]{figures/channel_comparison.pdf}
\caption{Mean LOSO Pearson correlation for each (aligner, channel set)
combination on the 23 subjects. Error bars are standard deviation
across subjects. Numeric labels above each bar report the mean COR.
Significance markers above each channel set compare Riemannian vs.\
None (${}^{***}p<0.001$; n.s.\ = not significant). The red annotation
indicates that the 4-channel forehead setup retains 98.3\% of the
17-channel Riemannian performance.}
\label{fig:channels}
\end{figure}

\subsection{Ablations}
\label{sec:ablation}

The following ablation studies are conducted on a 10-subject subset to
verify the contributions of individual pipeline components. The 10 subjects
were selected as the first 10 subjects in the dataset's standard ordering
(i.e., sub01--sub10), which is a fixed, protocol-independent ordering that
does not involve any outcome-based selection. All ablations
use the 17-channel setup with Riemannian alignment unless otherwise noted.

\subsubsection{Reference point} We compare aligning each subject to the
identity (Riemannian mean of its own SPD matrices) with aligning to the
grand mean of all training subjects. The two reference points yield
comparable regression performance on a 10-subject subset and on the
full 23-subject dataset, consistent with the theoretical expectation
that the Riemannian geometry is invariant to the choice of reference.
We use the identity (per-subject mean) as the default reference
throughout this work, as it avoids the additional cost of estimating
a global Riemannian mean.

\subsubsection{Covariance estimator} We compare OAS, LWF, and SCM under
both None and Riemannian alignment on 10 subjects. Without alignment,
OAS yields the best COR ($0.596$), followed by LWF ($0.589$) and SCM
($0.516$). With Riemannian alignment, all estimators improve: OAS
reaches $0.635$, LWF reaches $0.628$, and SCM reaches $0.581$. The
conditioning of the SPD matrices matters: SCM produces matrices with
larger eigenvalue spread, which propagates into the tangent space
projection. OAS is therefore the default estimator.

\subsubsection{Regressor} We compare RidgeCV and linear-kernel SVR under
both None and Riemannian alignment. Without alignment, RidgeCV and
SVR are similar ($0.596$ vs $0.585$). With Riemannian alignment, both
regressors improve substantially (RidgeCV $0.635$, SVR $0.638$),
confirming that the alignment benefit is regressor-agnostic and
compatible with both regularized linear and large-margin models.

\subsubsection{Per-band vs.\ global alignment} Aligning each band
independently (the proposed default) reaches COR $0.635$ on the
10-subject ablation, while aligning a single global mean across all
bands reaches COR $0.631$. The per-band default is preferable because
different frequency bands exhibit different inter-subject variability
patterns, although the difference is modest ($+0.004$).

\subsection{Frequency-band importance}

Fig.~\ref{fig:band} shows the distribution of the selected
features across the five frequency bands, aggregated over all
LOSO folds of the Riemannian-aligned pipeline. The $\alpha$ band
contributes the largest share (28.6\%, 429 features), closely followed by
$\beta$ (21.0\%, 315 features) and $\theta$ (22.7\%, 341 features);
together $\alpha$ and $\beta$ account for
approximately 50\% of the selected features. The remaining bands
contribute $\gamma$ (17.3\%, 259 features) and $\delta$
(10.4\%, 156 features). This distribution is consistent across all three
aligners (None, Euclidean, Riemannian), indicating that the frequency
band importance is driven by the underlying EEG physiology rather than
the alignment method. The dominance of $\alpha$ and $\theta$ bands
aligns with the physiological literature on vigilance, which
associates decreases in $\alpha$ power with drowsiness and increases
in $\theta$ activity during reduced alertness.

\begin{figure}[!t]
\centering
\includegraphics[width=\columnwidth]{figures/band_importance.pdf}
\caption{Distribution of the \texttt{SelectKBest}-selected features
across the five EEG frequency bands under Riemannian alignment,
aggregated over all LOSO folds. Each bar
is labeled with the percentage of total selected features and the
absolute count. The $\alpha$, $\theta$, and $\beta$ bands together account for
approximately 72\% of the selected features, consistent with
vigilance-related EEG patterns.}
\label{fig:band}
\end{figure}

\subsection{Comparison with Riemannian transfer-learning methods}

Table~\ref{tab:sota} compares the parallel transport alignment with two
Riemannian transfer-learning baselines---RPA
\cite{Rodrigues2019} and the Zanini et al.\ alignment
\cite{zanini2018}---under the same LOSO protocol on the 17-channel
setup. RPA applies parallel transport followed by a scaling factor
$\lambda$; the Zanini alignment transports each subject to the grand
mean of all training subjects rather than to the identity.

On the 17-channel setup, RPA reaches COR $0.672 \pm 0.196$, which
is above the Euclidean aligner ($0.655$) but below the proposed
method ($0.692$). The Zanini alignment yields comparable mean COR to
the proposed method ($0.692$), which is expected because both methods
use Riemannian parallel transport and differ only in the reference
point (identity vs.\ grand mean); this similarity is discussed in
Section~III-D (ablation) and the Methods section. On the 4-channel forehead setup, the proposed method
(COR $0.681$) outperforms both RPA ($0.650$) and the Euclidean
aligner ($0.642$), with the Wilcoxon signed-rank test confirming
significance over baseline ($p < 0.0001$, $d = 0.850$).
The Zanini alignment performs similarly on the forehead setup
($0.681$), confirming that the choice of reference point does not
affect regression performance on reduced channel configurations.
We use the identity reference as the default throughout this work
for its simplicity and computational efficiency.

\begin{table}[!t]
\caption{Head-to-head comparison of Riemannian transfer-learning
methods on the 17-channel LOSO setup (23 subjects).}
\label{tab:sota}
\centering
\setlength{\tabcolsep}{4pt}
\begin{tabular}{l c c}
\toprule
\textbf{Method} & \textbf{COR} & \textbf{RMSE} \\
\midrule
None (baseline)   & $0.585 \pm 0.230$ & $0.278$ \\
EA \cite{he2018euclidean} & $0.655 \pm 0.218$ & $0.224$ \\
RPA \cite{Rodrigues2019}  & $0.672 \pm 0.196$ & $0.200$ \\
Zanini et al.\ \cite{zanini2018} & $0.692 \pm 0.206$ & $0.195$ \\
\textbf{Ours}     & $\mathbf{0.692} \pm \mathbf{0.206}$ & $\mathbf{0.195}$ \\
\bottomrule
\end{tabular}
\end{table}

\subsection{Statistical robustness}

To complement the paired $t$-test, we report non-parametric and
resampling-based statistics. A Wilcoxon signed-rank test on the
17-channel setup yields $p = 7.9 \times 10^{-6}$, and on the
4-channel forehead setup $p = 8.8 \times 10^{-5}$, both well below
the $\alpha = 0.05$ threshold. Bootstrap resampling ($n = 2000$)
gives a 95\% confidence interval for Cohen's $d$ of $[0.62, 1.21]$
on the 17-channel setup and $[0.60, 1.31]$ on the 4-channel setup;
neither interval crosses zero, confirming that the effect is robust
to the normality assumption. A permutation test with $10{,}000$
random label flips yields $p < 0.0001$ for the 17-channel setup and
$p = 0.0001$ for the 4-channel setup, providing distribution-free
confirmation of significance.

We note that the main analysis involves 3 channel configurations
$\times$ 3 alignment methods = 9 paired comparisons, plus 4 ablation
comparisons. With a Bonferroni correction for 13 comparisons, the
adjusted $\alpha$ threshold is $0.05/13 \approx 0.0038$. The primary
result---Riemannian vs.\ None on the 17-channel setup---yields a
paired $t$-test $p = 0.0007$, which remains significant after
Bonferroni correction. Similarly, the 4-channel forehead comparison
(Riemannian vs.\ None, $p = 0.0001$) passes the correction. We
therefore report uncorrected $p$-values but note that the main
conclusions are unchanged after correction.

\subsection{Computational cost}

Table~\ref{tab:compcost} reports the wall-clock time for each
pipeline stage on a single subject (885 epochs, 17 channels). The
filter-bank and covariance estimation stage (7.2~s) is identical for
both aligners because it is applied before alignment. The
Riemannian mean computation itself adds negligible overhead
($<0.01$~s per subject) relative to the no-alignment baseline. The
full LOSO fold takes approximately $64$~s per subject under Riemannian
alignment vs.\ $54$~s without alignment. The dominant overhead comes from
feature selection and regression training, while alignment increases
the total fold time by approximately $10$~s (18.6\%), almost entirely
from the feature selection and regression stage ($43.6$~s vs.\
$53.6$~s). This increase is likely due to the Riemannian-aligned
features having a larger eigenvalue spread in the concatenated
tangent space, which increases the condition number of the design
matrix and slows the convergence of the RidgeCV internal
leave-one-out cross-validation grid search. The alignment step
itself, in contrast, adds negligible overhead ($<0.01$~s per subject).
This overhead is acceptable for offline analysis;
for real-time wearable deployment, the pre-trained model eliminates
the per-fold cost entirely, leaving only the filter-bank and
covariance-estimation stages.

\begin{table}[!t]
\caption{Per-subject wall-clock time (seconds) for each pipeline
stage, averaged over 5 subjects (17 channels, 885 epochs).}
\label{tab:compcost}
\centering
\begin{tabular}{l c c}
\toprule
\textbf{Stage} & \textbf{None} & \textbf{Riemannian} \\
\midrule
Filter + covariance       & 7.2 & 7.2 \\
Riemannian mean           & 1.58 & 1.57 \\
Tangent space             & 1.74 & 1.73 \\
Feature sel.\ + regression & 43.6 & 53.6 \\
\midrule
Full LOSO fold            & 54.1 & 64.2 \\
\bottomrule
\end{tabular}
\end{table}

\section{Discussion}
\label{sec:discussion}

\subsubsection{Why Riemannian geometry outperforms Euclidean alignment} The persistent gap between Euclidean and Riemannian
alignment (Table~\ref{tab:main}) is consistent with the theoretical
argument that the arithmetic mean is not a natural summary on a
curved space. The Euclidean mean of SPD matrices is biased toward the
cone apex when the matrices are heterogeneous, and whitening by the
Euclidean mean propagates that bias into the tangent space. The
Riemannian mean is invariant to the parameterization of the manifold
and is the geometrically correct center. The empirical gap between the
two aligners supports the view that the curved-space notion of center
matters in practice for EEG covariance features. A rival explanation
is that the gain is driven by the per-band normalization alone; the
ablations argue against this, because the per-band normalization is
common to both aligners and the Euclidean aligner does not reach
Riemannian performance under the same normalization. These findings
are consistent with the review by Yger et al. \cite{yger2017riemannian},
which concluded that Riemannian geometry provides the largest benefit
when the SPD matrices are heterogeneous, precisely the condition that
characterises cross-subject EEG data.

\subsubsection{Per-subject alignment is sufficient at test time} Our protocol aligns
the test subject using only its own SPD matrices, with no leakage
from the training subjects. This is a departure from the BCI
literature in which the test subject is sometimes left unaligned. We
find that unaligned test subjects degrade the LOSO correlation
substantially (consistent with the findings of
\cite{zanini2018}); aligning the test subject using its own mean is
both necessary and sufficient.

\subsubsection{Compatibility with downstream refinements} The benefit is preserved when the EEG feature set is
reduced from 17 channels to 4 forehead channels (98.3\% of full
Riemannian performance), and when the covariance estimator is swapped
for LWF. The benefit is also preserved across the RidgeCV and linear
SVR regressors, indicating that the gain comes from the alignment step
rather than from a particular regressor.

The 6-channel temporal configuration showed no significant improvement
over baseline (Table~\ref{tab:main}), which merits further discussion.
The temporal baseline COR ($0.659$) is already markedly higher than
the 17-channel baseline ($0.585$), indicating that the inherent
cross-subject variability in the temporal region is substantially
lower. This pattern is physiologically plausible: the temporal EEG
sites (T7, T8, etc.) lie directly over the lateral cortex and are less
affected by variations in skull thickness, scalp impedance, and muscle
artifacts that disproportionately impact frontal and parietal
electrodes. Consequently, the subject-specific SPD shift that Riemannian
alignment is designed to correct is naturally smaller in the temporal
region, leaving limited room for alignment-driven improvement. This
finding aligns with the broader BCI observation that Riemannian
geometry provides the largest gains when inter-subject heterogeneity
is high \cite{yger2017riemannian}. Notably, Euclidean alignment on the
6-channel temporal setup actually decreases COR relative to no
alignment ($0.622$ vs.\ $0.659$), indicating that the larger
eigenvalue spread of the temporal SPD matrices makes the arithmetic
mean a poor whitening reference, which introduces distortion rather
than correction.

\subsubsection{Relation to existing Riemannian transfer-learning methods}
Table~\ref{tab:sota} shows that the proposed method yields comparable
performance to the Zanini et al.\ alignment \cite{zanini2018} on the
17-channel setup, which is expected because both use Riemannian
parallel transport and differ only in the reference point (identity
vs.\ grand mean); the ablation confirms that this choice has minimal
impact on regression performance.
We emphasize that the principal contribution of this work is not the
alignment mechanism itself, but the first systematic demonstration that
Riemannian geometry-based alignment is both necessary and effective for
\emph{continuous vigilance regression}---a fundamentally different
problem from the classification tasks on which Zanini et al.\ and
subsequent works were validated. Vigilance is a high-frequency,
continuously varying signal whose dynamics impose stricter demands on
the preservation of SPD manifold structure. Our experiments further
provide the first evidence that Riemannian parallel transport alignment
translates robustly to a 4-channel wearable EEG configuration, a
practical scenario that has not been examined in prior Riemannian BCI
transfer-learning studies. RPA
\cite{Rodrigues2019} underperforms the proposed method by $0.020$
COR on the 17-channel setup, likely because the scaling factor
$\lambda$ introduces an additional degree of freedom that is
difficult to estimate reliably with the limited calibration data
available in the LOSO setting. Zhuo et al.\ \cite{zhuo2024} proposed
log-Euclidean Riemannian transfer learning for EEG classification,
achieving comparable performance to affine-invariant methods at lower
computational cost; our affine-invariant approach remains tractable
for the 17-channel vigilance regression task (Table~\ref{tab:compcost}).
Ju and Guan \cite{ju2025} combined optimal transport with SPD manifold
geometry for cross-session BCI, which is complementary to our
parallel-transport alignment and could potentially be integrated in
future work. Paillard et al.\ \cite{paillard2025} recently showed that
integrating Riemannian geometry with learnable spectral filters
achieves state-of-the-art performance across multiple EEG prediction
tasks, further supporting the value of geometry-aware representations.

\subsubsection{Comparison with deep domain adaptation methods}
Recent years have seen a surge of deep learning-based domain adaptation
approaches for EEG, including domain-adversarial neural networks
(DANN), conditional domain adversarial networks (CDAN), and
Wasserstein-distance alignment. These methods are powerful for
high-data regimes but face several practical limitations in the
vigilance regression setting. First, deep domain adaptation models
typically require thousands of labeled samples per subject for
effective training, whereas the LOSO setting provides only 22
training subjects with modest per-subject sample sizes. Second, the
computational demands of deep models are incompatible with real-time
wearable deployment: as shown in Table~\ref{tab:compcost}, the
proposed Riemannian pipeline incurs a total overhead of approximately
$10$~s per fold, with the alignment step itself adding less than
$0.01$~s per subject, making it suitable for on-device deployment
even on a 4-channel wearable EEG system. Third, deep domain
adaptation often requires careful hyperparameter tuning and large
batch sizes to avoid mode collapse or overfitting, whereas the
proposed Riemannian approach is deterministic and parameter-free
beyond the reference point choice. We therefore regard the
Riemannian geometry framework as a pragmatic and well-grounded
alternative to deep methods for small-sample cross-subject
regression, particularly in resource-constrained wearable settings.

\subsubsection{Practical implications} The main LOSO result (Table
\ref{tab:main}) achieves COR $0.692$ with $p=0.0007$ and Cohen's
$d=0.824$, demonstrating that Riemannian parallel transport alignment
is both statistically significant and practically meaningful for
cross-subject generalization. The 4-channel forehead configuration
retains 98.3\% of the full 17-channel performance, making the approach
viable for wearable EEG systems with limited channel coverage.
The computational overhead of Riemannian alignment is modest
(Table~\ref{tab:compcost}): the Riemannian mean computation adds
less than $0.01$~s per subject, and the full LOSO fold overhead is
approximately $10$~s, dominated by the filter-bank and
covariance-estimation stages rather than the alignment itself.

\subsubsection{Limitations} (i) The main LOSO result is on a single dataset (SEED-VIG). Future work
will explore cross-dataset generalization and shared alignment models,
including validation on publicly available fatigue and
affective-EEG datasets such as SADT \cite{zhao2021lammsda} and
DEAP \cite{koelstra2012deap}, to assess the robustness of Riemannian
alignment across diverse EEG paradigms and recording conditions.
(ii) The covariance matrices are computed from 8-second epochs;
shorter or longer windows may yield different behaviour. (iii) The
alignment step is fit per subject; an open question is whether a
single shared alignment can be learned across datasets. (iv) The
current pipeline uses fixed hyperparameters; an adaptive scheme that
selects the calibration set size per subject is left to future work.

\section{Conclusion}
\label{sec:conclusion}

We presented a Riemannian parallel transport alignment as a
pre-projection step for cross-subject EEG vigilance regression. On
SEED-VIG with 23 subjects under LOSO, alignment increases the mean
Pearson correlation from $0.585$ to $0.692$ ($p = 0.0007$, $d =
0.824$) and improves 21 of 23 subjects individually. The benefit
is statistically significant on both the full 17-channel setup and
the 4-channel forehead setup, retaining 98.3\% of the
full-setup performance. A head-to-head comparison with RPA and the
Zanini et al.\ alignment confirms that the proposed method yields
comparable performance to existing Riemannian transfer-learning
baselines, and
non-parametric tests (Wilcoxon, permutation) and bootstrap confidence
intervals corroborate the significance of the improvement. Ablations
isolate the gain to the Riemannian geometry rather than individual
pipeline choices. Future work will explore shared alignment models
across datasets and adaptive per-subject calibration sizing. The
result is restricted to short-epoch EEG covariance features on a
single dataset; its transfer to long-window features, other vigilance
proxies, or cross-dataset settings remains to be established.

\subsubsection{Data and Code Availability}
The SEED-VIG dataset used in this study is publicly released by the
SJTU BCMI Laboratory for non-commercial research use and is
available to qualified researchers on request from the dataset
administrator at \protect\url{https://bcmi.sjtu.edu.cn/home/seed/seed-vig.html}
\cite{zheng2017}. The dataset is a third-party resource and is not
redistributed by the authors.

\section*{Acknowledgments}
This work was supported by the Shenzhen Science and Technology
Program under Grant KCXFZ20240903094011015. The authors thank the
SJTU SEED Lab for sharing the SEED-VIG dataset.

\begin{thebibliography}{22}

\bibitem{gao2024}
D. Gao, P. Li \emph{et~al.}, ``CSF-GTNet: A novel multi-dimensional
feature fusion network based on Convnext-GeLU-BiLSTM for
EEG-signals-enabled fatigue driving detection,'' \emph{IEEE J.
Biomed. Health Inform.}, vol. 28, no. 5, pp. 2558--2568, May 2024.

\bibitem{cui2023}
J. Cui \emph{et~al.}, ``Benchmarking EEG-based cross-dataset driver
drowsiness recognition with deep transfer learning,'' in \emph{Proc.
45th Annu. Int. Conf. IEEE Eng. Med. Biol. Soc. (EMBC)}, 2023,
pp.~1--5.

\bibitem{zheng2017}
W.-L. Zheng and B.-L. Lu, ``A multimodal approach to estimating
vigilance using EEG and forehead EOG,'' \emph{J. Neural Eng.},
vol. 14, no. 2, 026017, Feb. 2017.

\bibitem{zanini2018}
P. Zanini, M. Congedo, C. Jutten, S. Said, and Y. Berthoumieu,
``Transfer learning: A Riemannian geometry framework with
applications to brain-computer interfaces,'' \emph{IEEE Trans.
Biomed. Eng.}, vol. 65, no. 5, pp. 1107--1116, May 2018.

\bibitem{Rodrigues2019}
P. L. C. Rodrigues, C. Jutten, and M. Congedo, ``Riemannian
Procrustes analysis: Transfer learning for brain-computer
interfaces,'' \emph{IEEE Trans. Biomed. Eng.}, vol. 66, no. 8,
pp. 2390--2401, Aug. 2019.

\bibitem{lotte2018review}
F. Lotte, L. Bougrain, A. Cichocki, M. Clerc, M. Congedo,
A. Rakotomamonjy, and F. Yger, ``A review of classification
algorithms for EEG-based brain-computer interfaces: A 10 year
update,'' \emph{J. Neural Eng.}, vol. 15, no. 3, 031005,
Jun. 2018.

\bibitem{moakher2005}
M. Moakher, ``A differential geometric approach to the geometric
mean of symmetric positive-definite matrices,'' \emph{SIAM J. Matrix
Anal. Appl.}, vol. 26, no. 3, pp. 735--747, 2005.

\bibitem{wu2022tnsre}
D. Wu, Y. Xu, and B.-L. Lu, ``Transfer learning for EEG-based
brain-computer interfaces: A review of progress made since 2016,''
\emph{IEEE Trans. Cogn. Dev. Syst.}, vol. 14, no. 1, pp. 4--19,
Mar. 2022.

\bibitem{he2018euclidean}
H. He and D. Wu, ``Transfer learning for brain-computer interfaces:
A Euclidean space data alignment approach,'' \emph{IEEE Trans.
Biomed. Eng.}, vol. 67, no. 2, pp. 399--410, Feb. 2020.

\bibitem{barachant2010lvaica}
A. Barachant, S. Bonnet, M. Congedo, and C. Jutten, ``Riemannian
geometry applied to BCI classification,'' in \emph{Proc. 9th Int.
Conf. Latent Variable Anal. Signal Separation (LVA/ICA)}, 2010,
pp. 629--636.

\bibitem{barachant2012tbme}
A. Barachant, S. Bonnet, M. Congedo, and C. Jutten, ``Multiclass
brain-computer interface classification by Riemannian geometry,''
\emph{IEEE Trans. Biomed. Eng.}, vol. 59, no. 4, pp. 920--928,
Apr. 2012.

\bibitem{zhuo2024}
F. Zhuo, X. Zhang, F. Tang, Y. Yu, and L. Liu, ``Riemannian
transfer learning based on log-Euclidean metric for EEG
classification,'' \emph{Front. Neurosci.}, vol. 18, 1381572,
May 2024.

\bibitem{yair2019ptspd}
O. Yair, M. Ben-Chen, and R. Talmon, ``Parallel transport on the
cone manifold of SPD matrices for domain adaptation,''
\emph{IEEE Trans. Signal Process.}, vol. 67, no. 7, pp. 1797--1811,
Apr. 2019.

\bibitem{congedo2017review}
M. Congedo, A. Barachant, and R. Bhatia, ``Riemannian geometry for
EEG-based brain-computer interfaces; a primer and a review,''
\emph{Brain-Comput. Interfaces}, vol. 4, no. 3, pp. 155--174,
2017.

\bibitem{ju2025}
C. Ju and C. Guan, ``Deep optimal transport for domain adaptation
on SPD manifolds,'' \emph{Artif. Intell.}, vol. 345, 104347, 2025.

\bibitem{ang2012fbcsp}
K. K. Ang, Z. Y. Chin, C. Wang, C. Guan, and H. Zhang, ``Filter bank
common spatial pattern algorithm on BCI competition IV Datasets 2a
and 2b,'' \emph{Front. Neurosci.}, vol. 6, no. 39, pp. 1--10,
Mar. 2012.

\bibitem{chen2010oas}
Y. Chen, A. Wiesel, Y. C. Eldar, and A. O. Hero, ``Shrinkage
algorithms for MMSE covariance estimation,'' \emph{IEEE Trans.
Signal Process.}, vol. 58, no. 10, pp. 5016--5029, Oct. 2010.

\bibitem{yger2017riemannian}
F. Yger, M. Berar, and F. Lotte, ``Riemannian approaches in
brain-computer interfaces: A review,'' \emph{IEEE Trans. Neural
Syst. Rehabil. Eng.}, vol. 25, no. 10, pp. 1753--1762, Oct. 2017.

\bibitem{pennec2006}
X. Pennec, P. Fillard, and N. Ayache, ``A Riemannian framework
for tensor computing,'' \emph{Int. J. Comput. Vis.}, vol. 66,
no. 1, pp. 41--66, Jan. 2006.

\bibitem{paillard2025}
J. Paillard, J. F. Hipp, and D. A. Engemann, ``GREEN: A
lightweight architecture using learnable wavelets and Riemannian
geometry for biomarker exploration with EEG signals,''
\emph{Patterns}, vol. 6, no. 3, 101182, Mar. 2025.

\bibitem{zhao2021lammsda}
Y. Zhao, G. Dai, G. Borghini, J. Zhang, X. Li, Z. Zhang,
P. Aric\`{o}, G. Di Flumeri, F. Babiloni, and H. Zeng,
``Label-based alignment multi-source domain adaptation for
cross-subject EEG fatigue mental state evaluation,''
\emph{Front. Hum. Neurosci.}, vol. 15, 706270, Oct. 2021.

\bibitem{koelstra2012deap}
S. Koelstra, C. M\"{u}hl, M. Soleymani, J.-S. Lee, A. Yazdani,
T. Ebrahimi, T. Pun, A. Nijholt, and I. Patras, ``DEAP: A
database for emotion analysis using physiological signals,''
\emph{IEEE Trans. Affect. Comput.}, vol. 3, no. 1, pp. 18--31,
Jan.--Mar. 2012.

\end{thebibliography}

% \begin{IEEEbiography}{Huang Tan}
% received the B.S. degree in software engineering from
% Sichuan University, Chengdu, China, in 2021. He is currently
% pursuing the Ph.D. degree with the University of Electronic
% Science and Technology of China, Chengdu. His research
% interests include biomedical signal processing, Riemannian
% geometry for brain-computer interfaces, and transfer learning
% for EEG-based vigilance estimation.
% \end{IEEEbiography}

% \begin{IEEEbiography}{Xiangzhu Li}
% is with the Department of Neurosurgery, Shenzhen Hospital of
% Integrated Traditional Chinese and Western Medicine, Shenzhen,
% China, where he is currently an Attending Physician. His
% clinical interests include neurosurgical treatment and
% EEG-based neurological monitoring.
% \end{IEEEbiography}

% \begin{IEEEbiography}{Li Zhang}
% received the M.D. degree from Fudan University Shanghai Medical
% College, Shanghai, China. He is currently a Chief Physician,
% Professor, and Head of Neurosurgery at the Shenzhen
% Hospital of Integrated Traditional Chinese and Western
% Medicine, Shenzhen, China. His clinical expertise includes
% microsurgical treatment of complex cerebral aneurysms,
% arteriovenous malformations, arteriovenous fistulas, carotid
% endarterectomy, brain tumors, trigeminal neuralgia, and
% hemifacial spasm; neuroendoscopic minimally invasive
% techniques for cranial diseases; multidisciplinary management
% of severe cerebrovascular disease in elderly patients; and
% cross-disciplinary treatment of severe traumatic brain injury
% and hypertensive intracerebral hemorrhage.
% \end{IEEEbiography}

% \begin{IEEEbiography}{Guangqiang Yin}
% is a Professor with the School of Information and Software
% Engineering and the Director of the Public Security Technology
% Research Center, University of Electronic Science and
% Technology of China (UESTC), Chengdu, China. He is also a
% National-Level Leading Talent under the Shenzhen High-Level
% Talent Program. His research interests include intelligent
% image processing, information security, and public safety
% technologies.
% \end{IEEEbiography}

\end{document}


